\originalchapter{约束、阻力与连续绳}\label{ch:constraint}
\originalsection{绳长约束与滑轮}
连接多物体时, 各物体加速度并不一定相等. 不可伸长绳给出几何约束, 再对时间求导得到速度和加速度约束. 例如固定滑轮两端竖直向下坐标为 $y_1,y_2$, 则 $y_1+y_2=\text{常数}$, 所以 $a_1=-a_2$. 质量 $m_2>m_1$ 的 Atwood 装置中,
\begin{equation}
m_2g-T=m_2a,\qquad T-m_1g=m_1a,
\end{equation}
故 $a=(m_2-m_1)g/(m_1+m_2)$, $T=2m_1m_2g/(m_1+m_2)$. 有惯量的滑轮必须另写力矩方程, 见第\ref{ch:rotation}章.

若动滑轮由两段同绳承托, 它向下移动 $y$ 会同时增加两段长度, 自由端坐标 $x$ 满足 $x+2y=\text{常数}$. 因而自由端速度为动滑轮速度的反向两倍; 对动滑轮受力时两段张力相加为 $2T$. 动滑轮的机械优势来自这组几何关系, 不是凭直觉分配张力.

一般约束 $f(q_1,\ldots,q_n,t)=0$ 的求导须包含显含时间项, 例如移动滑轮. 运动方程与约束一起求解后还要核对接触、摩擦及绳张力的可实现条件.

\originalsection{线性及平方阻力}
阻力模型取决于速度、介质与物体形状. 小速率某些黏性流动可用 $\vct F_d=-b\vct v$, $b$ 单位 kg/s. 竖直向下为正, 从静止下落满足 $m\dot v=mg-bv$. 把 $v_t=mg/b$, $\tau=m/b$ 代入, 对 $v-v_t$ 分离变量, 得
\begin{equation}
v=v_t(1-e^{-t/\tau}),\qquad
z=v_t\bigl[t-\tau(1-e^{-t/\tau})\bigr].
\end{equation}
$t\ll\tau$ 时展开指数恢复 $v\simeq gt$, $z\simeq gt^2/2$; $t\gg\tau$ 时速度趋于终端速度 $v_t$, 合力趋于零而重力并未消失.

较高速率常用 $\vct F_d=-c|\vct v|\vct v$, $c$ 单位 kg/m. 向下运动时 $m\dot v=mg-cv^2$, 终端速度 $v_t=\sqrt{mg/c}$. 积分 $\mathrm dv/(1-v^2/v_t^2)=g\dd t$ 给
\begin{equation}
v=v_t\tanh(gt/v_t),\qquad
z=\frac{v_t^2}{g}\ln\cosh(gt/v_t).
\end{equation}
这组解限于从静止下落、阻力系数恒定; 上升段阻力符号改变. 量纲估算 $c\sim\rho R^2$ 仅表示几何面积与密度的依赖, 阻力系数还可依赖形状与无量纲流动参数.

\originalsection{有质量绳与缠绕摩擦}
对竖直均匀绳, 自由端到截面的距离为 $s$, 线密度 $\lambda$, 静止时截面需承托下方质量, 所以 $T(s)=\lambda gs$. 若整绳以加速度 $a$ 向上, 下方绳段满足 $T-\lambda sg=\lambda sa$, 得 $T=\lambda s(g+a)$. 非均匀绳只需把下方质量换成 $\int\lambda(s')\dd s'$. 从底端边界 $T(0)=0$ 出发比猜测顶端力更可靠.

绳绕圆柱的接触角为 $\Theta$, 摩擦系数 $\mu$. 对转角 $\dd\theta$ 的无质量绳段, 径向平衡的一阶形式给法向力 $\dd N=T\dd\theta$, 切向平衡给 $|\dd T|\le\mu T\dd\theta$. 若即将向紧端滑动, 等号成立, $\mathrm dT/T=\mu\dd\theta$, 积分得到
\begin{equation}
T_{\rm tight}/T_{\rm loose}=e^{\mu\Theta}.
\end{equation}
未达到滑动极限时只有不等式 $T_{\rm tight}/T_{\rm loose}\le e^{\mu_s\Theta}$. 绕行角必须用弧度; 多绕一圈会乘上 $e^{2\pi\mu}$, 这说明很小的牵引端张力也能约束很大负载.
\originalsection{微段方程与边界}
令竖直绳自自由底端向上坐标为 $s$, 局部线密度 $\lambda(s)$. 取长 $\dd s$ 的小段, 上端向上张力 $T(s+\dd s)$, 下端向下张力 $T(s)$, 重量 $\lambda(s)g\dd s$. 若向上加速度为 $a$, 则
\begin{equation}
T(s+\dd s)-T(s)-\lambda(s)g\dd s=\lambda(s)a\dd s.
\end{equation}
除以 $\dd s$ 取极限, 得 $\mathrm dT/\mathrm ds=\lambda(s)(g+a)$, 积分 $T(s)=T(0)+(g+a)\int_0^s\lambda(u)\dd u$. 若下端悬质量 $m_0$, 边界变成 $T(0)=m_0(g+a)$; 若自由则为零. 微分方程不给边界常数, 边界来自末端实际载荷.
